% %%%%%%%%%%%%%%%%%%%%%%%%%%%%%%%%%%%%%%%%%%%%%%%%%%%%%%%%%%%%%%%%%%%%%%%%%%%%%%%%%
%
%	DADOS.tex 
%  -----------
%
%	Neste arquivo você escreve os dados necessários do seu trabalho.
%
% %%%%%%%%%%%%%%%%%%%%%%%%%%%%%%%%%%%%%%%%%%%%%%%%%%%%%%%%%%%%%%%%%%%%%%%%%%%%%%%%%
%
%
% Titulo, local, data e data completa
\titulo{Ferramenta Computacional para o Dimensionamento de Vigas em Concreto Armado com Abertura}
\local{Porto Alegre}
\data{2022}
\qualdatacompleta{9 de dezembro de 2026} % ou [Mounth] [day^{th}] 202X}
%
% Palavras-chave
\quaispalavraschave{vigas, concreto armado, abertura, método da rigidez, estado limite último}
\quaiskeywords{beam, reinforced concrete, openings, stiffness method, ultimate limit state}
\quaispalabrasclave{vigas; hormigón armado; aberturas; método de la rigidez; estado límite último}
%
% Dados do autor
\autor{Cícero Mangili Nicolini}
\qualcitarautor{Nicolini, C. M.}
\qualemailautor{cicero.augusto12@gmail.com}
%
% Dados do orientador
\orientador{Prof. Felipe Pinto da Motta Quevedo}
\qualtitorientador{Dr. pela Universidade Federal do Rio Grande do Sul}
\qualgeneroorientador{M}
%
% Dados coorientador
\temcoorientador{Não}
\coorientador{Profa. Enedina Alves Marques}
\qualtitcoorientador{Eng. pela Universidade Federal do Paraná}
\qualgenerocoorientador{F}
%
% Dados do coordenador do curso
\qualcoordenador{Nilo Consoli}
\qualtitcoordenador{Ph.D. pela Concordia University, Canadá}
\qualgenerocoordenador{M}
%
% Dados sobre a banca (no máximo 6)
\quantosmembrosbanca{3}
%
% Adicionando membros da banca e dados deles
\addtext{Prof. Luís Carlos Prestes}					% Nome da banca
\addtext{UFRGS}										% Instituição da banca
\addtext{Ph.D. pela Universidade de Origem, País}	% Titulação da banca
%
\addtext{Profa. Elmina Tessa Wilson}
\addtext{ISU}
\addtext{Eng. pela Iowa State University}
%
\addtext{Prof. Leonel de Moura Brizola}
\addtext{UFRGS}
\addtext{Dr. pela~\imprimirinstituicao}
%
\addtext{Profa. Olive Wetzel Dennis}
\addtext{Cornell University}
\addtext{Md pela Comlumbia University}
%